\documentstyle[%


righttag%


,11pt%


,amssymb%


,russian%


,izvru%


]{amsart}





\newcommand{\tl}[3]{\medskip\noindent{\bf#1}\ #2 \dotfill #3 \par} 


\pagestyle{empty}


% \pagestyle{plain}


\begin{document}


\par


\vskip 2cm


\par


\bigskip





\bigskip


\par


\centerline{\Large СОДЕРЖАНИЕ}


\bigskip


\bigskip


%%%%%%%%%% Use \newline %%%%%%%%%%%








\tl{Бобочко В.Н.}


{Дифференциальная точка поворота


в теории сингулярных возмущений. I}{3}


\tl{Кац Б.А., Погодина А.Ю.}


{О граничных значениях интеграла типа Коши на негладкой\newline


кривой}{15}


\tl{Кокурин М.Ю.}


{Условие истокопредставимости и оценки скорости 


сходимости методов\newline регуляризации линейных уравнений в 


банаховом пространстве. II}{22}


\tl{Корешков Н.А.} 


{Элементы Казимира $\Bbb Z$-форм модулярных алгебр Ли}{32}


\tl{Марюкова Н.Е.}


{Поверхности постоянной кривизны


в квазипсевдоримановом пространстве\newline


постоянной кривизны и уравнение Клейна--Гордона}{36}


\tl{Тронин С.Н.}


{Операды в категории конвексоров. I}{42}


\tl{Фернандеш В.У.}


{Новый класс делителей полугрупп изотонных отображений


конечных\newline цепей}{51}


\tl{Шишкин Г.И.}


{Оптимальные по порядку скорости сходимости кусочно-равномерные


сетки\newline для сингулярно возмущенных уравнений конвекции--диффузии}{60}








\bigskip


\par


\centerline{\Large КРАТКИЕ~СООБЩЕНИЯ}


\bigskip





\tl{Галимянов А.Ф.}


{К прямым методам решения интегральных уравнений 


с\newline логарифмически ослабленными ядрами Коши 


на разомкнутых контурах}{73}


\tl{Эскин Л.Д.}


{Об интегральном уравнении, описывающем ориентационные


фазовые\newline переходы в модели Парсонса для системы осесимметричных частиц}{78}











\vfill


\noindent\raisebox{2pt}{\copyright} Казанский госудаpственный унивеpситет


\end{document}


